% Generated draft: requires mathematical and source review.
\documentclass[11pt]{article}
\usepackage[margin=25mm]{geometry}
\usepackage{fontspec}
\setmainfont{Latin Modern Roman}
\usepackage{amsmath,amssymb,mathtools,unicode-math}
\setmathfont{Latin Modern Math}
\usepackage{enumitem,needspace,xurl,hyperref}
\setlength{\parindent}{0pt}
\setlength{\parskip}{5pt}
\setlength{\emergencystretch}{3em}
\newcommand{\N}{\mathbb{N}}
\newcommand{\Z}{\mathbb{Z}}
\newcommand{\Q}{\mathbb{Q}}
\newcommand{\R}{\mathbb{R}}
\newcommand{\C}{\mathbb{C}}
\newcommand{\F}{\mathbb{F}}
\newcommand{\A}{\mathbb{A}}
\newcommand{\PP}{\mathbb{P}}
\newcommand{\E}{\mathbb{E}}
\newcommand{\eps}{\varepsilon}
\newcommand{\dd}{\,\mathrm{d}}
\begin{document}
\section*{UnsolvedMath ID: 30002814}
\textbf{Source problem code:} catalogue entry
\subsection{Definitions and mathematical statement}
Let \(\Omega\subset\C\) be a Jordan domain, meaning the bounded domain enclosed by a Jordan curve, and let \(p,q,r\in\partial\Omega\) be pairwise distinct. For each mesh size \(\varepsilon>0\), the source prescribes a maximal simply connected region \(\Omega_\varepsilon\subset\Omega\) triangulated by equilateral Euclidean triangles of side length \(\varepsilon\), with boundary vertices \(p_\varepsilon,q_\varepsilon,r_\varepsilon\) satisfying
\[
 |p_\varepsilon-p|\leq 2\varepsilon,\qquad
 |q_\varepsilon-q|\leq 2\varepsilon,\qquad
 |r_\varepsilon-r|\leq 2\varepsilon.
\]
Here the distances are Euclidean. The admissible class relative to which the region is maximal is not specified in the supplied record.

For a triangulated region with disk topology, its metric double is obtained from two isometric copies by identifying corresponding boundary points, equipped with the induced intrinsic length metric. Denote the resulting doubled surface by \(S_\varepsilon\), and denote the images of the three selected boundary vertices by \(P_\varepsilon,Q_\varepsilon,R_\varepsilon\). On a Euclidean polyhedral surface, the curvature at a vertex \(v\) is its angle defect
\[
 K(v)=2\pi-\sum_{f\ni v}\theta_f(v),
\]
where \(\theta_f(v)\) is the angle of an incident triangular face at \(v\).

The source asks for a discrete conformal uniformization of \(S_\varepsilon\) to a Euclidean polyhedral sphere with metric \(\widehat d_\varepsilon\) such that
\[
 \operatorname{Area}(S_\varepsilon,\widehat d_\varepsilon)=\frac32,
 \qquad
 K(P_\varepsilon)=K(Q_\varepsilon)=K(R_\varepsilon)=\frac{4\pi}{3},
\]
with zero curvature at every other vertex. The supplied record does not define the discrete conformal equivalence relation or specify the existence and uniqueness hypotheses for this uniformization; these remain part of the unresolved formulation.

\textbf{Conjecture as supplied, with unresolved qualifications.} For every Jordan domain and every triple of distinct boundary points as above, do the resulting uniformized polyhedral spheres converge, as \(\varepsilon\to0\), to the metric double of the unit equilateral triangle?

The source does not specify the notion of convergence, whether markings or identifications with the original domain are retained, or whether the assertion must hold for every admissible choice of triangulations and boundary vertices. Its prescription for every \(\varepsilon>0\) also lacks an existence qualification for the approximating regions; a formulation restricted to sufficiently small admissible mesh sizes would require clarification.

There is an additional normalization issue: if ``unit equilateral triangle'' means side length \(1\), its metric double has Euclidean area \(\sqrt{3}/2\), whereas the prescribed approximating spheres have area \(3/2\). The intended meaning of ``unit,'' the area normalization, and the convergence notion must therefore be clarified together. Both source normalizations are retained here without correction.
\subsection{Short English statement}
The supplied conjecture starts with equilateral triangulations inside a Jordan domain, doubles each triangulated region across its boundary, and applies discrete conformal uniformization so that curvature is concentrated equally at three selected boundary points. It asks whether these spheres converge to a doubled equilateral triangle. The meaning of maximal approximation, discrete conformality, convergence, and the target scale remains unresolved. The source labels the problem open; this draft does not independently verify that status.
\subsection{Sources}
\begin{itemize}
\item Supplied local UnsolvedMath registry record, permanent ID \(30002814\), title: \emph{Discrete Riemann Mapping by Conformal Triangulations}; source problem code \texttt{catalogue entry}. Collection: Oberwolfach Reports Open Problems. The record gives proposed year \(2015\) and does not name an individual proposer.
\item Supplied report citation: \emph{Discrete Differential Geometry} (2015), Oberwolfach Reports, DOI 10.4171/owr/2015/13, pp.\ 661--729. \url{https://doi.org/10.4171/owr/2015/13}. The background describes Oberwolfach Reports as documenting mathematical workshops and their open-problem discussions. No problem-specific page or passage is supplied.
\item Supplied literature source: Bobenko--Pinkall--Springborn. \url{https://arxiv.org/abs/1005.2698}. The record supplies this attribution and link but no title or specific supporting passage.
\item Literature checked at: \texttt{2026-08-21}. This is supplied dated curation, not independently verified. The supplied assessment labels the conjecture open and reports that no proof or counterexample was found in the checked primary literature, listing OWR Report 13/2015 and Bobenko--Pinkall--Springborn. The background explicitly calls this dated literature triage derived from the supplied OWR curation and requests independent verification of time-sensitive status and cited claims.
\item The registry labels the statement \texttt{corrected\_verified} and reports a manual semantic, typographic, and LaTeX review. These are supplied metadata claims, not independent mathematical or source verification by this draft.
\end{itemize}
\end{document}
